\unnumberedchapter{Introduction}
\epigraph{\textcolor{Cerulean}{\bf{Logic:}} Do it for the picture puzzle. Put it all together. Solve the world. One conversation at a time.}{Disco Elysium}



As scientists, and physicists in particular, we discover Brownian motion as early as in high school, since it is perfectly well understood and predictable for the study of simple cases. Typically, we are given an aqueous solution containing small particles (colloids). We observe them under an optical microscope and marvel, because our scientific culture is not yet fully developed, at the unstoppable, unpredictable motion of the colloids, even though the fluid surrounding them is perfectly quiescent. We then learn that what we are observing is none other than one of the first experimental proofs of the atomic nature of matter. Our instrument is restained in its observational capabilities, and its lower size limit sits several orders of magnitude beyond the size of a molecule. What we actually see is the consequence of the existence of these atoms, rather than the atoms themselves: the temperature of the fluid translates into the motion of molecules at their own scale. These molecules, though minute compared to the size of a colloid, move at high speeds: several hundred meters per second on average. By striking a colloid, they transfer a bit of their momentum with each impact, and since these impacts are themselves random both in time and direction, the colloid is set into an equally unpredictable, erratic motion.

All these observations have been extensively confirmed and formalised, and I describe the historical case of simple Brownian motion in the first chapter of this manuscript. I also lay out in this chapter the various theoretical developments important for a proper understanding of my doctoral work, and refer back to them throughout the manuscript. Moreover, Brownian motion is not merely a textbook example, constructed to make us infer molecular reality. It is in fact an extremely common process in many biological systems. Going back to its first historical observation in pollen grains, we already note that the shape of the colloids itself leads to modifications of the canonical theories; this in fact constitutes one of the earliest generalisations brought to Brownian motion theory. One may also wonder, since this motion arises from thermal agitation within the fluid, what happens when the fluid itself is disturbed, for example by modifying its properties (viscosity, temperature), or by leaving the dilute limit where colloids are treated as infinitely far apart, entering a dense regime where the fluid displaced by the motion of one colloid directly affects that of another, giving rise to hydrodynamic coupling. Another natural line of exploration concerns the effect of rigid walls, which likewise modify the behavior of the fluid and therefore of the colloid's motion. These questions have been explored, and partly resolved, over the past century, but their complexity does not always allow for a complete and full understanding of their unification. I explore this in chapter 2 and provide some insights. Then comes the matter of the influence of boundary deformability in confined systems. Living matter is full of deformable membranes: a cell wall, a mucous membrane, a soft tissue. Sometimes these membranes are themselves subject to the same motion arising from collisions with fluid molecules. How can we predict the influence of their deformability and their fluctuations on a neighboring colloid? Could we control, enhance, or reduce the mobility of a molecule for therapeutic purposes? Conversely, might it be possible to halt the spread of a pathogen across such membranes, a pathogen responsible for the sometimes incurable degeneration of certain of our capacities? This problem is inherently complex to solve, but in chapter 3 I describe an elementary theoretical model to qualitatively capture the effect of deformability on Brownian motion. The flow of the fluid itself also affects the motion of colloids, first and most obviously because the colloid gets carried along by the flow, but also for less trivial reasons that I develop in chapter 4. In addition, I seek to understand whether the density, even low, of a colloidal solution in a flow can affect the dynamics of the colloids themselves. Finally, we have already touched on many scenarios specific to the living world, but what about living matter itself? Colloids can, in some cases, move on their own by using the energy available around them. When they are too large to be visibly affected by the thermal agitation of the fluid, we can nonetheless predict their behavior by treating them, by extension, as effective Brownian particles. It is then possible to extract further information from their trajectories, just as we would for a small passive colloid. In chapter 5, I use this type of model to study the motion of microswimmers of a size much larger than that of a Brownian particle, in particular to study the behavior of these microswimmers in a flow.


\section*{Open-access}
All of the results I produced throughout my thesis rely on numerical simulation as the main tool. I chose, from the start of my thesis, to approach research from this angle, since it permits different effects to be properly isolated in order to better predict their consequences. It would not be possible, experimentally, to observe a perfect sphere, an absolute zero temperature, a strictly parabolic flow, or infinitely pure water. For me, each of these is at most a few lines of code. Naturally, with this capability comes the question of how ``real'' what I simulate actually is. I must always keep a highly skeptical eye on my results, and compare them against theory and experiment whenever these exist. I see numerical tools as a fantastic Swiss Army knife, allowing both the exploration of effects that are experimentally inaccessible due to current technological limitations, and the validation of theories that are unrealisable because of the laws of nature, even though such theories still contribute, if not quantitatively then at least conceptually, vital developments for research.
Since I trained throughout my thesis using freely accessible resources (forums, publications, tutorials) as well as human interactions, the two pillars of research, all the code I developed is publicly available on \href{https://github.com/Laacherez/thesis}{dedicated GitHub repository}\footnote{https://github.com/Laacherez/thesis/Codes}. I place immense importance on open science, and I believe that a researcher's success should not be limited to the funding they are granted, but should rest on their curiosity, their ability to ask the right questions, to inform themselves and collaborate to answer them, and on their pedagogy and integrity in communicating their results.